\section{Ziel}
Vor diesem Hintergrund zielt das Projekt darauf ab, ein interaktives Visualisierungs-Framework zu entwickeln, das die Lehre im Bereich der Robotik an der Fachhochschule Dortmund unterstützt und bereichert. Das Framework wird speziell für den Einsatz innerhalb der JupyterLab-Umgebung konzipiert und soll dazu dienen, kinematische Zusammenhänge auf intuitive und anschauliche Weise darzustellen.

Ein zentrales Anliegen ist es, Lehrpersonen und Aufgabenstellern ein Werkzeug an die Hand zu geben, mit dem sie durch einfache Python-Befehle individualisierbare Aufgabenstellungen und Übungsszenarien gestalten können. Diese sollen es Studierenden ermöglichen, eigenständig mit Parametern zu experimentieren, Bewegungsabläufe zu analysieren und dadurch ein tieferes Verständnis für die Wirkungsweise kinematischer Modelle zu entwickeln.

Besonderes Augenmerk liegt auf der Darstellung verschiedener Rotationsmodelle im Raum, darunter Eulerwinkel, Roll-Pitch-Yaw-Winkel, Rotationsmatrizen und Quaternionen. Durch deren visuelle Umsetzung werden nicht nur Unterschiede und Gemeinsamkeiten der jeweiligen Darstellungsformen sichtbar, sondern auch deren konkrete Bedeutung im räumlichen Kontext erfahrbar gemacht.

Insgesamt verfolgt das Projekt das Ziel, den Lernprozess im Bereich der Kinematik durch moderne, interaktive Lehrmittel zu unterstützen und einen Beitrag zu einer zeitgemäßen, verständnisorientierten Ausbildung in der Robotik zu leisten.